Notes:

- Modulus:
\[
2197=13^3.
\]
- Need \(m^{2028}\equiv1\pmod {13^3}\) and no earlier exponent \(>1\) gives \(1\).  So I use
\[
\operatorname{ord}_{13^3}(m)=2028.
\]

The group of units has size
\[
\phi(13^3)=13^2\cdot 12=2028.
\]
It is cyclic.  The number of generators should be
\[
\phi(2028).
\]
Using \(2028=2^2\cdot 3\cdot 13^2\), I calculate
\[
\phi(2028)=\phi(2^2)\phi(3)\phi(13^2)=2\cdot2\cdot13=52.
\]

Now count lifts to \(0,\ldots,9999\).  There are four full periods modulo \(2197\), with an extra initial segment:
\[
9999=4\cdot2197+1211.
\]
So each good residue occurs four times, and good residues in the first initial part occur one extra time.  Let \(n_5\) be the number of good residues in that first piece.  Then
\[
\#=4\cdot 52+n_5=208+n_5.
\]

I estimate \(n_5\) by uniform distribution:
\[
n_5\approx 52\cdot\frac{1210}{2196}\approx 28.6.
\]
Taking \(n_5=28\), I get
\[
\#=236,\qquad P=\frac{236}{10000}=\frac{59}{2500}.
\]
Therefore
\[
p+q=59+2500=2559\equiv \boxed{559}\pmod {1000}.
\]
